% \section{Estimating the Rademacher and Gaussian Complexities of Function Classes}
% \section{Estimating the Complexities}

% \begin{frame}{Structural Results}

% Various estimations with respect to
% \begin{itemize}
%     \item Convex hull.
%     \item Lipschitz composition (notation: $\phi\circ F = \{\phi\circ f\,|\, f\in F\}$).
%     \item Minkowski sum.
% \end{itemize}
% We will mention when needed later.
% \end{frame}

% \begin{frame}{Motivation}
% \textbf{Main question:} how can we estimate the complexity of a complicated class $F$?

% \medskip
% Even though the empirical Rademacher complexity can be estimated from one sample and one realization of the signs, computing it still requires solving
% \begin{align*}
%     \sup_{f\in F}\sum_{i=1}^n \sigma_i f(X_i),
% \end{align*}

% which is difficult for rich classes. Instead, many large classes can be built from simpler basis classes:
% \begin{itemize}
%     \item decision trees = Boolean combinations of node functions,
%     \item voting methods = thresholded convex combinations,
%     \item neural networks = Lipschitz activations composed with linear combinations.
% \end{itemize}

% \medskip
% Hence, we try to establish the \emph{structural bounds}:
% \begin{align*}
%     \text{complexity of }F
% \quad \Longrightarrow \quad
% \text{complexities of simpler classes }F_1,\dots,F_m.
% \end{align*}

% \end{frame}

% \begin{frame}{Empirical estimation: Theorem 11}
%     \begin{mytheorem}
%         Let $F$ be a class of functions mapping to $[-1,1]$. Then for any integer $n$,
%         \begin{align*}
%             \mathbb{P}\!\left(
% \left|
% R_n(F)-\frac{2}{n}\sup_{f\in F}\left|\sum_{i=1}^n \sigma_i f(X_i)\right|
% \right|\ge \varepsilon
% \right)
% \le 2\exp\!\left(-\frac{\varepsilon^2 n}{8}\right),
%         \end{align*}
%         and 
%         \begin{align*}
%             \mathbb{P}\!\left(
% \left|R_n(F)-\widehat R_n(F)\right|\ge \varepsilon
% \right)
% \le 2\exp\!\left(-\frac{\varepsilon^2 n}{8}\right).
%         \end{align*}
%     \end{mytheorem}

%     \begin{itemize}
%     \item $R_n(F)$ is well approximated by the empirical quantity $\widehat R_n(F)$;
%     \item so Rademacher complexity is genuinely \emph{data-dependent} and \emph{estimable}.
% \end{itemize}
% However, computing the supremum over $F$ is still hard. That motivates the structural results in the rest of the section.
% \end{frame}

% \begin{frame}{Some simple structural results: Theorem 12}

%     \begin{mytheorem}
%     Let $F, F_1,..., F_k$ and $H$ be classes of real functions. Then
%         \begin{enumerate}
%     \item If $F\subseteq H$, then $ R_n(F)\le R_n(H).$

%     \item $R_n(F)=R_n(\operatorname{conv}F)=R_n(\operatorname{absconv}F).$

%     \item For every $c\in\mathbb{R}$, $R_n(cF)=|c|\,R_n(F).$

%     \item If $\phi:\mathbb{R}\to\mathbb{R}$ is Lipschitz with constant $L_\phi$ and $\phi(0)=0$, then
%     \[
%     R_n(\phi\circ F)\le 2L_\phi\,R_n(F).
%     \]

%     \item For any uniformly bounded $h$,
%     \[
%     R_n(F+h)\le R_n(F)+\frac{\|h\|_\infty}{\sqrt n}.
%     \]

    
% \end{enumerate}
%     \end{mytheorem}
% \end{frame}

% \begin{frame}{Some simple structural results: Theorem 12}
%     \begin{mytheorem}
%         \begin{enumerate}
%         \setcounter{enumi}{5}
%             \item For $1\le q<\infty$, let
%     \[
%     L_{F,h,q}:=\{|f-h|^q:\ f\in F\}.
%     \]
%     If $\|f-h\|_\infty\le 1$ for every $f\in F$, then
%     \[
%     R_n(L_{F,h,q})\le 2q\left(R_n(F)+\frac{\|h\|_\infty}{\sqrt n}\right).
%     \]

%     \item We have
%     \[
%     R_n\!\left(\sum_{i=1}^k F_i\right)\le \sum_{i=1}^k R_n(F_i).
%     \]
%         \end{enumerate}
%     \end{mytheorem}
% \end{frame}

% \subsection{Lipschitz function on $\mathbb R^k$}

% \begin{frame}{Lipschitz functions on $\mathbb{R}^k$}
%     Theorem 12(4) handled real-valued composition:
%     \begin{align*}
%         f \mapsto \phi\circ f,\qquad \phi:\mathbb{R}\to\mathbb{R}.
%     \end{align*}
%     Now we need a vector-valued version. Suppose $F \subseteq F_1 \oplus \cdots \oplus F_m,$ then each $f\in F$ has the form $f(x) = (f_1(x),\dots,f_m(x)), f_i\in F_i.$

%     Let
% \[
% \phi:Y\times \mathbb{R}^m \to \mathbb{R}
% \]
% be Lipschitz in the second argument. 

% \textbf{Question:} how large can the Gaussian complexity of $\phi\circ F$ be?
% \end{frame}

% \begin{frame}{Theorem 14}
%     \begin{alertblock}{Lemma}
%         Let $\{X_i\}_{i=1}^m$ and $\{Y_i\}_{i=1}^m$ be Gaussian processes such that
% \[
% \|X_i-X_j\|_2^2 \le \|Y_i-Y_j\|_2^2
% \qquad \text{for all } i,j.
% \]
% Then
% \[
% \mathbb{E}\sup_i X_i \le 2\,\mathbb{E}\sup_i Y_i.
% \]
%     \end{alertblock}

% Using the bound, assume also that for each $y\in Y$, the map
% \[
% \phi(y,\cdot):\mathbb{R}^m\to \mathbb{R}
% \]
% is Lipschitz with constant $L$, passes through the origin, and is uniformly bounded.

% Then for every sample $(X_1,Y_1),\dots,(X_n,Y_n)$,
% \[
% \widehat G_n(\phi\circ F)
% \le
% 2L\sum_{i=1}^m \widehat G_n(F_i).
% \]
    
% \end{frame}




% \subsection{Boolean Combinations of Functions}

% \begin{frame}{3.3 Boolean combinations of functions}
% Let
% \[
% g:\{\pm1\}^k \to \{\pm1\}
% \]
% be a fixed Boolean function, and let $F_1,\dots,F_k$ be classes of $\{\pm1\}$-valued functions.

% Then
% \[
% G_n(g(F_1,\dots,F_k))
% \le
% 2\sum_{j=1}^k G_n(F_j).
% \]

% \medskip
% \textbf{Why this matters:}
% many classifiers are built by combining simple binary rules:
% \begin{itemize}
%     \item AND / OR / majority,
%     \item tree node combinations,
%     \item logical decision architectures.
% \end{itemize}

% So Boolean composition does not destroy complexity control.
% \end{frame}